\chapter{Deliverable}
\label{chap:deliverable}

\section{Software Solution}
\label{section:software-solution}

Vid-Scratch is delivered as a desktop application built with Tauri,
providing a self-contained adversarial video poisoning pipeline that
runs entirely on the user's local machine. The source code is available
at the project GitHub repository.

This section describes the full poisoning pipeline from video input to
adversarial video output, followed by screenshots of the application
user interface.

% ─────────────────────────────────────────────────────────────────────
\subsection{Poisoning Pipeline}
\label{subsection:poisoning-pipeline}

The pipeline transforms an input video into an adversarially perturbed
video through six sequential stages, as illustrated below.

\begin{figure}[H]
    \centering
    % Replace with your pipeline diagram if available
    \includegraphics[width=\textwidth]{chapter_7/pipeline.png}
    \caption{Overview of the Vid-Scratch adversarial poisoning pipeline.}
    \label{fig:pipeline-overview}
\end{figure}

% ── Stage 1 ──────────────────────────────────────────────────────────
\subsubsection{Stage 1: Video Loading and Frame Sampling}

The input video is decoded frame by frame using OpenCV. All
$N$ frames are loaded at their original resolution
$(H_{\text{orig}} \times W_{\text{orig}})$.
A model window of $T = 40$ uniformly sampled frames
is then selected from the full sequence using linear indexing:

\begin{equation}
    i_t = \left\lfloor
        \frac{t \cdot (N - 1)}{T - 1}
    \right\rfloor,
    \quad t = 0, 1, \ldots, T-1
    \label{eq:frame-sampling}
\end{equation}

\noindent
where $i_t$ is the index of the $t$-th sampled frame in the
original sequence. The sampled frames are then resized to
$224 \times 224$ pixels for model input.

% ── Stage 2 ──────────────────────────────────────────────────────────
\subsubsection{Stage 2: Poison Frame Budget}

The number of frames to poison, $k$, is computed automatically
from the total frame count:

\begin{equation}
    k = \max\!\left(1,\ \mathrm{round}\!\left(\frac{N}{40}\right)\right)
    \label{eq:poison-budget}
\end{equation}

\noindent
For a short video of $N < 40$ frames, $k = \max(1, \lfloor 0.25 N \rfloor)$
is used instead to ensure at least one frame is poisoned.
The budget $k$ controls how many frames in the 40-frame model
window will receive adversarial perturbations.

% ── Stage 3 ──────────────────────────────────────────────────────────
\subsubsection{Stage 3: Bayesian Frame Selection (Algorithm 1)}

Rather than perturbing all $T$ frames equally,
Vid-Scratch selects the $k$ \emph{most vulnerable} frames
using Bayesian Optimisation (BO) with a Gaussian Process (GP)
surrogate model.

A candidate binary mask $\mathbf{m} \in \{0,1\}^T$ assigns
value~1 to $k$ selected frames and~0 to the rest.
For each candidate mask, a proxy adversarial loss is evaluated:

\begin{equation}
    \mathcal{L}_{\text{proxy}}(\mathbf{m}) =
        \lambda \, \mathcal{L}_{\text{SSIM}}(\hat{X}, X)
        - \mathcal{L}_{\text{CE}}(f(\hat{X}),\, y)
    \label{eq:proxy-loss}
\end{equation}

\noindent
where $X$ is the clean model-window tensor,
$\hat{X}$ is its perturbed version under mask $\mathbf{m}$,
$f(\cdot)$ is the SlowFast R101 model,
$y$ is the ground-truth label,
$\mathcal{L}_{\text{CE}}$ is the cross-entropy loss, and
$\mathcal{L}_{\text{SSIM}}$ is the SSIM-based perceptual
similarity loss.

The GP models the mapping from mask to proxy loss.
At each BO iteration, the next candidate is chosen by
maximising the Expected Improvement (EI) acquisition function:

\begin{equation}
    \mathrm{EI}(\mathbf{m}) =
        \mathbb{E}\!\left[
            \max\!\left(
                \hat{\mathcal{L}}(\mathbf{m}) - \mathcal{L}^{*} - \xi,\; 0
            \right)
        \right]
    \label{eq:ei}
\end{equation}

\noindent
where $\mathcal{L}^{*}$ is the best observed proxy loss so far
and $\xi > 0$ is an exploration parameter.
After $N_{\text{BO}}$ iterations (default: 50),
the mask with the highest proxy loss is selected.

% ── Stage 4 ──────────────────────────────────────────────────────────
\subsubsection{Stage 4: Adversarial Generation (Algorithm 2)}

Given the selected mask $\mathbf{m}^*$, the adversarial perturbation
is optimised jointly over a spatial flow field $U \in \mathbb{R}^{T \times 2 \times H \times W}$
and an additive noise tensor $N \in \mathbb{R}^{T \times C \times H \times W}$.

The perturbed frame sequence is constructed as:

\begin{equation}
    \hat{X}_t =
    \begin{cases}
        \left( \mathcal{ST}(X_t,\, U_t) + N_t \right)_{[0,1]} & \text{if } m_t = 1 \\
        X_t & \text{if } m_t = 0
    \end{cases}
    \label{eq:perturbed-frame}
\end{equation}

\noindent
where $\mathcal{ST}(\cdot, U_t)$ is a differentiable spatial
transformer that warps frame $X_t$ according to flow field $U_t$,
and $(\cdot)_{[0,1]}$ denotes pixel-level clamping to $[0,1]$.

The full objective minimised by Adam is:

\begin{equation}
    \mathcal{L} =
        \lambda \, \mathcal{L}_{\text{SSIM}}\!\left(
            \hat{X}_{\mathcal{S}},\; X_{\mathcal{S}}
        \right)
        - \mathcal{L}_{\text{CE}}\!\left(
            f(\hat{X}),\; y
        \right)
    \label{eq:main-loss}
\end{equation}

\noindent
where $\mathcal{S} = \{t : m_t = 1\}$ is the set of poisoned frames.
Minimising $\mathcal{L}$ simultaneously encourages high
perceptual similarity on poisoned frames (via the SSIM term)
and maximises the classification loss so that the model is
more likely to misclassify (via the adversarial term).

At each optimisation step, the noise is clamped:

\begin{equation}
    N_t \leftarrow \mathrm{clip}\!\left(N_t,\; -\epsilon,\; +\epsilon\right)
    \label{eq:noise-clamp}
\end{equation}

\noindent
where $\epsilon$ is the noise budget (default: $\epsilon = 0.03$).
Optimisation also terminates early if the SSIM value on poisoned
frames drops below the SSIM budget threshold $1 - \delta_{\text{SSIM}}$
(default: $\delta_{\text{SSIM}} = 0.03$).

% ── Stage 5 ──────────────────────────────────────────────────────────
\subsubsection{Stage 5: Error-Minimising Noise (Training-Time Layer)}

For the training-time attack scenario, a second noise layer
is applied to the \emph{non-attacked} frames using the
Error-Minimising (EM) noise strategy. This layer is designed
to reduce the usefulness of the poisoned video as a training
sample by minimising the classification loss on the remaining
frames.

The EM noise $\delta_t$ is updated iteratively using
momentum-based gradient descent with sign gradient:

\begin{equation}
    \mu_{s+1} = \beta \, \mu_s
        + \frac{\nabla_{\delta} \mathcal{L}_{\text{CE}}\!\left(
            f(X_t + \delta_t),\, y
        \right)}{\left\lVert
            \nabla_{\delta} \mathcal{L}_{\text{CE}}
        \right\rVert_1 + \epsilon_0}
    \label{eq:em-momentum}
\end{equation}

\begin{equation}
    \delta_{t} \leftarrow \delta_{t} - \alpha \cdot \mathrm{sign}(\mu_{s+1})
    \label{eq:em-update}
\end{equation}

\noindent
where $\beta = 0.9$ is the momentum coefficient,
$\alpha$ is the step size, and $\epsilon_0$ is a small
stabiliser. The noise is constrained by
$\lVert \delta_t \rVert_\infty \leq \epsilon_{\text{EM}}$
(default: $\epsilon_{\text{EM}} = 0.002$),
keeping the base layer imperceptible to human viewers.

% ── Stage 6 ──────────────────────────────────────────────────────────
\subsubsection{Stage 6: Perturbation Mapping and Output}

The adversarial perturbation is computed at the 224-pixel
model resolution. To produce the output video at the original
resolution, the difference $\Delta_t = \hat{X}_t^{(224)} - X_t^{(224)}$
is upscaled back:

\begin{equation}
    \Delta_t^{(\text{orig})} =
        \gamma \cdot \mathrm{Resize}\!\left(
            \Delta_t,\; H_{\text{orig}} \times W_{\text{orig}}
        \right)
    \label{eq:perturbation-mapping}
\end{equation}

\noindent
where $\gamma$ is a codec boost factor (default: $\gamma = 1.5$)
that compensates for perturbation attenuation caused by video
compression. The final output frame is:

\begin{equation}
    X_t^{(\text{out})} =
        \mathrm{clip}\!\left(
            X_t^{(\text{orig})} + \Delta_t^{(\text{orig})},\; 0,\; 1
        \right)
    \label{eq:output-frame}
\end{equation}

\noindent
All $N$ original frames are preserved in the output video,
with only the $k$ poisoned frames carrying the adversarial
perturbation. The output video is saved at the original
resolution, frame rate, and with audio from the source file
remuxed without re-encoding.

% ─────────────────────────────────────────────────────────────────────
\subsection{User Interface}
\label{subsection:user-interface}

The Vid-Scratch desktop application provides a three-step
workflow: select a video, configure settings, and run the
poisoning pipeline. The following figures illustrate the
main screens of the application.

\begin{figure}[H]
    \centering
    \includegraphics[width=0.85\textwidth]{chapter_7/UI_Open.png}
    \caption{
        Initial state of the application.
        The left panel shows the three-step workflow.
        The right panel shows two video preview areas
        (Original and Poisoned), both empty before a
        video is selected.
    }
    \label{fig:ui-open}
\end{figure}

\begin{figure}[H]
    \centering
    \includegraphics[width=0.85\textwidth]{chapter_7/UI_Put_Video.png}
    \caption{
        After the user selects an input video via the
        native file dialog, the original video is loaded
        into the preview panel and the selected file path
        is displayed. The Poisoned panel remains empty
        until the pipeline has been executed.
    }
    \label{fig:ui-put-video}
\end{figure}

\begin{figure}[H]
    \centering
    \includegraphics[width=0.85\textwidth]{chapter_7/UI_Start_Process.png}
    \caption{
        Processing state. After the user clicks
        \textit{Start Process}, the button changes to
        \textit{Processing...} and a progress bar appears
        below the output destination showing the current
        stage (e.g., ``Loading model... 5\%'').
        The Poisoned panel displays a \textit{Processing...}
        placeholder while the pipeline runs.
    }
    \label{fig:ui-processing}
\end{figure}

\begin{figure}[H]
    \centering
    \includegraphics[width=0.85\textwidth]{chapter_7/UI_Finish.png}
    \caption{
        Completed state. Both the original and poisoned
        videos are shown side by side in the preview panels.
        The output section displays a
        \textit{Protection successful!} message along with
        the total processing time and number of retry
        attempts.
        A \textit{Show Predictions} button appears to
        reveal the model prediction panel.
    }
    \label{fig:ui-finish}
\end{figure}

\begin{figure}[H]
    \centering
    \includegraphics[width=0.85\textwidth]{chapter_7/UI_Prediction.png}
    \caption{
        Prediction panel expanded.
        The right side shows two prediction blocks:
        \textit{Clean — Top-5 Predictions} displays the
        model's confidence scores on the original video
        (e.g., \textit{squat} at 100.0\%),
        and \textit{Poisoned — Top-5 Predictions} displays
        the model's predictions on the poisoned video
        (e.g., \textit{eating watermelon} at 39.9\%),
        confirming that the attack successfully caused
        the model to misclassify the action.
    }
    \label{fig:ui-prediction}
\end{figure}

\section{Test Report}
\label{section:test-report}

The testing of Vid-Scratch is divided into three parts:
an inference-time benchmark on real Kinetics-400 videos,
a training-time poisoning experiment, and a unit test
suite covering the core pipeline modules.

% ─────────────────────────────────────────────────────────────────────
\subsection{Inference-Time Attack Benchmark}
\label{subsection:inference-benchmark}

The inference-time attack was evaluated on 10 video clips
sampled from the Kinetics-400 dataset using the default
pipeline configuration~\cite{kay2017kinetics}
($\epsilon = 0.03$, $\delta_{\text{SSIM}} = 0.03$,
$N_{\text{BO}} = 50$, $T_{\text{max}} = 150$).
Each video was processed by the full VidScratch pipeline
and the output was verified by reloading the saved file
and running SlowFast R101 inference again to confirm
that the attack holds after codec compression.

The results are summarised in Table~\ref{tab:benchmark}.

\begin{table}[H]
\centering
\caption{Inference-time attack results on 10 Kinetics-400 videos.
Orig = original predicted class index,
Adv = adversarial predicted class index,
SSIM = structural similarity between original and poisoned video.}
\label{tab:benchmark}
\renewcommand{\arraystretch}{1.2}
\begin{tabular}{clcccc}
\hline
\textbf{\#} & \textbf{Video} & \textbf{Orig} & \textbf{Adv} & \textbf{SSIM} & \textbf{Time (s)} \\
\hline
1  & 01R6cmvFbZI & 108 & 353 & 0.9441 & 47.6 \\
2  & -2VKVjgNuE0 & 305 & 105 & 0.9555 & 35.0 \\
3  & --SrV9bGzRA & 310 & 160 & 0.9548 & 51.0 \\
4  & 0G2A0x2-vf4 & 117 & 114 & 0.9497 & 35.1 \\
5  & -ayNZ5bs7EQ & 289 & 256 & 0.9237 & 34.9 \\
6  & -9GW\_m6MQd0 & 103 &   1 & 0.9268 & 45.1 \\
7  & -7mJ12n-xyM & 335 & 304 & 0.9726 & 42.1 \\
8  & -3zcsBnDVzU & 145 & 363 & 0.9542 & 41.4 \\
9  & 0fjsXnccCu0 & 318 & 134 & 0.9768 & 46.4 \\
10 & -2FbowxhzfQ & 301 & 222 & 0.9428 & 41.6 \\
\hline
\multicolumn{2}{l}{\textbf{Attack success rate}} & \multicolumn{4}{r}{\textbf{10/10 (100\%)}} \\
\multicolumn{2}{l}{\textbf{Verified success rate}} & \multicolumn{4}{r}{\textbf{10/10 (100\%)}} \\
\multicolumn{2}{l}{\textbf{Average SSIM}} & \multicolumn{4}{r}{\textbf{0.9501}} \\
\multicolumn{2}{l}{\textbf{Average time}} & \multicolumn{4}{r}{\textbf{42.0 s/video}} \\
\hline
\end{tabular}
\end{table}

All 10 videos were successfully misclassified both in memory
and after saving and reloading the output file, giving a
verified attack success rate of 100\%. The average SSIM
across all poisoned videos was 0.9501, indicating that the
perturbations remain visually close to the originals.
Average processing time was 42.0 seconds per video on CPU.

% ─────────────────────────────────────────────────────────────────────
\subsection{Training-Time Poisoning Experiment}
\label{subsection:training-experiment}

To evaluate the training-time attack scenario, a small
classification model was trained under three data conditions
and evaluated on a held-out clean test set at every epoch
over 80 epochs. The three conditions were:

\begin{enumerate}
    \item \textbf{Clean model:} trained on 100\% clean videos.
    \item \textbf{Poison model:} trained on 100\% poisoned videos.
    \item \textbf{Mixed model:} trained on 50\% poisoned and 50\% clean videos.
\end{enumerate}

The dataset consisted of two action classes (lunges and squats)
drawn from a UCF-101 subset~\cite{soomro2012ucf101}, with 60 videos per class and a
50/60 train split. The target model was a lightweight
two-class classifier fine-tuned from SlowFast R101 features.
Training used a batch size of 4, a learning rate of 0.0005,
and ran for 80 epochs.

\begin{figure}[H]
    \centering
    \includegraphics[width=\textwidth]{chapter_7/poison_comparison_3x1_loss.png}
    \caption{
        Training loss (white) and test loss on clean data (blue)
        across 80 epochs for three training conditions.
        Left: clean model. Centre: poison model.
        Right: mixed model (50\% poison + 50\% clean).
    }
    \label{fig:poison-loss}
\end{figure}

\begin{figure}[H]
    \centering
    \includegraphics[width=\textwidth]{chapter_7/poison_comparison_3x1_acc.png}
    \caption{
        Training accuracy (white) and test accuracy on clean data (blue)
        across 80 epochs for three training conditions.
        Left: clean model. Centre: poison model.
        Right: mixed model (50\% poison + 50\% clean).
    }
    \label{fig:poison-acc}
\end{figure}

The results show that all three models achieve high training
accuracy (above 90\%) by the end of training, which is
expected as each model is able to overfit its own training
distribution. The key observation is in the test loss on
clean data. The clean model shows a clear downward trend
in test loss, reflecting genuine generalisation to unseen
clean samples. The mixed model, by contrast, shows a
slightly increasing trend in test loss on clean data
despite the training loss continuing to decrease, which is
consistent with the poisoned samples introducing conflicting
gradient signals that partially disrupt clean generalisation.

These results should be interpreted with caution for two
reasons. First, the dataset is small (approximately 60
videos per class), which causes high variance in the loss
curves and makes it difficult to draw statistically strong
conclusions. Second, the poisoning was applied using a
SlowFast R101 model trained on all 400 Kinetics classes,
whereas the downstream classifier was trained on only two
classes, which reduces the alignment between the adversarial
perturbation and the training objective. Despite these
limitations, the mixed model's test loss trend suggests
that poisoned data can introduce measurable interference
in clean generalisation even at low noise levels, which
is consistent with the expected behaviour of the
error-minimising noise component.

% ─────────────────────────────────────────────────────────────────────
\subsection{Unit Testing}
\label{subsection:unit-testing}

Unit tests covering the core pipeline modules are
maintained in the project GitHub repository. The test
suite verifies the correctness of the video loading and
frame sampling functions in \texttt{video\_io.py}, the
SlowFast R101 wrapper in \texttt{model\_wrapper.py},
the Bayesian frame selector and adversarial generator
in \texttt{vidscratch.py}, and the perturbation mapping
and output saving functions.

All unit tests pass on the current codebase. No
continuous integration or deployment pipeline is
configured for this project, as the system is
distributed as a local desktop application and does
not involve a hosted backend service.
